\chapter{Fall lösen und Produkt}\label{ch:produkt}

\section{Den Fall Lena lösen}\label{sec:fall_loesen}

\begin{myexercise}[Die Auflösung]
	Tragen Sie Ihre drei Fallakten zusammen. Schreiben Sie Lena eine kurze Nachricht (4--6 Sätze):
	Was ist passiert, und welche \textbf{drei} Dinge sollte sie ab heute anders machen?
	\schreibzeile
	\schreibzeile
	\schreibzeile
	\schreibzeile
	\schreibzeile
	\begin{myanswer}
		Mögliche Punkte: Standort in der Kamera ausschalten bzw. Fotos ohne Ortsangaben verschicken;
		für jedes Konto ein eigenes, langes Passwort verwenden und die \gls{TFA} einschalten;
		bei Nachrichten mit Zeitdruck oder Passwortabfrage nie auf Links klicken;
		bei blossstellenden Bildern die Löschung verlangen, melden und Hilfe holen.
	\end{myanswer}
\end{myexercise}

\section{Ihr Produkt}\label{sec:produkt_waehlen}

Wählen Sie im Team \textbf{eines} der beiden Produkte. Zeit: etwa 30 Minuten.

\subsection{Produkt 1: Entscheidungsbaum \enquote{Darf ich das posten?}}

Gestalten Sie auf einem A3-Plakat einen Entscheidungsbaum, mit dem eine Mitschülerin in wenigen Sekunden prüfen kann, ob sie ein Bild oder Video posten darf.
Der folgende Anfang zeigt die Idee. Ergänzen Sie mindestens drei weitere Fragen (z.\,B. zu Lizenzen, Musik, Standort im Bild) und gestalten Sie das Plakat ansprechend.

\begin{center}
	\begin{tikzpicture}[
			node distance=0.9cm and 1.4cm,
			frage/.style={draw=RoyalBlue, fill=RoyalBlue!10, rounded corners, align=center, text width=4.2cm, minimum height=1cm},
			ja/.style={draw=ForestGreen, fill=ForestGreen!10, rounded corners, align=center, text width=3cm},
			nein/.style={draw=BrickRed, fill=BrickRed!10, rounded corners, align=center, text width=3cm},
			pfeil/.style={-{Stealth}, thick},
			lbl/.style={font=\small, fill=white, inner sep=1pt}
		]
		\node[frage] (f1) {Habe ich das Bild selbst gemacht?};
		\node[frage, below=of f1] (f2) {Sind Personen erkennbar?};
		\node[nein, right=of f1] (n1) {Lizenz prüfen \\ (\cref{ex:fallkarten}, Nr.~4 und 5)};
		\node[frage, below=of f2] (f3) {Sind alle einverstanden?};
		\node[ja, left=of f2] (j1) {\dots weitere Fragen \dots};
		\node[nein, right=of f3] (n2) {Nicht posten!};
		\node[ja, below=of f3] (j2) {\dots weitere Fragen \dots};

		\draw[pfeil] (f1) -- node[lbl] {nein} (n1);
		\draw[pfeil] (f1) -- node[lbl] {ja} (f2);
		\draw[pfeil] (f2) -- node[lbl] {nein} (j1);
		\draw[pfeil] (f2) -- node[lbl] {ja} (f3);
		\draw[pfeil] (f3) -- node[lbl] {nein} (n2);
		\draw[pfeil] (f3) -- node[lbl] {ja} (j2);
	\end{tikzpicture}
\end{center}

\subsection{Produkt 2: Phishing-Falle bauen}

Gestalten Sie auf einer Karte eine erfundene Phishing-Nachricht (SMS, E-Mail oder Chat), die möglichst echt wirkt und \textbf{mindestens drei} Warnsignale versteckt.
Notieren Sie die Lösung auf der Rückseite. Danach tauschen Sie die Karte mit einem anderen Team: Wer findet alle Warnsignale?

\begin{myattention}
	Verwenden Sie nur \textbf{erfundene} Firmen, Namen und Adressen. Keine echten Links, keine echten Personen, keine echten Logos.
\end{myattention}

\subsection{Checkliste für Ihr Produkt}

\begin{center}
	\renewcommand{\arraystretch}{1.4}
	\begin{tabular}{p{12.5cm}c}
		\textbf{Kriterium}                                                                   & \textbf{erfüllt}   \\
		\midrule
		Der Inhalt ist fachlich richtig (vergleichen Sie mit den Stationen A--C).            & \faSquare[regular] \\
		Das Produkt ist ohne Erklärung verständlich.                                         & \faSquare[regular] \\
		Es enthält mindestens drei eigene Ideen (Fragen bzw. Warnsignale).                   & \faSquare[regular] \\
		Es verletzt selbst keine Rechte (keine fremden Bilder ohne Lizenz, keine echten Personen). & \faSquare[regular] \\
	\end{tabular}
\end{center}

\section{Exit-Ticket}

\begin{myexercise}[Eine Sache ändere ich heute noch]
	Notieren Sie \textbf{eine} konkrete Sache, die Sie heute noch an Ihrem Handy oder Ihren Konten ändern werden.
	\schreibzeile
\end{myexercise}
